\documentclass[a4paper, final]{article}
\usepackage[russian]{babel}
\usepackage[utf8]{inputenc}
\usepackage[T1,T2A]{fontenc}
\usepackage[14pt]{extsizes}
\usepackage[left=20mm, top=20mm, right=20mm, bottom=20mm, footskip=10mm]{geometry}
\usepackage{ragged2e}
\usepackage{indentfirst}
\usepackage{listings}
\usepackage{minted}
\usepackage{hyperref}
\usepackage{graphicx}
\usepackage{tikz}
\usepackage{caption}
\usepackage{xcolor}
\usepackage{subcaption}
\usepackage{amssymb}
\usepackage{amsmath}
\usepackage{amssymb}
\usepackage{amsmath}
\usepackage{setspace}
\usepackage{relsize}
\definecolor{link_color}{HTML}{0000EE}

\usepackage{listings}
\lstset{
   language = C,
   basicstyle = \small \ttfamily,
   morecomment = [s][commentstyle]{/*}{*/},
   morecomment = [l][commentstyle]{//},
   commentstyle = \color{comment} \ttfamily \itshape,
   inputencoding = utf8,
   columns = fullflexible,
   showstringspaces = false,
   extendedchars = false
}

% Доопределение цвета, используемого в \lstset выше.
\definecolor{comment}{HTML}{6A9955}

\hypersetup{
    colorlinks=true,
    linkcolor=link_color,
    urlcolor=link_color,
    citecolor=link_color
}

\begin{document}

\thispagestyle{empty}
\begin{center}
    \hfill \break
    \hfill \break
    МИНИСТЕРСТВО НАУКИ И ВЫСШЕГО ОБРАЗОВАНИЯ РОССИЙСКОЙ ФЕДЕРАЦИИ\\[5 pt]
    Федеральное государственное автономное образовательное учреждение высшего образования
    «Санкт-Петербургский политехнический университет Петра Великого»\\[20 pt]
    Институт компьютерных наук и кибербезопасности\\ [10pt]
    Высшая школа технологий искусственного интеллекта\\ [10pt]
    Направление: 02.03.01 Математика и компьютерные науки\\[110pt]

    {\large «Цифровая аналитика»\\ [7.5 pt]
     Отчет о проделанной работе} \\ [7.5 pt]
    {\LARGE\bf «Разработка графического калькулятора с помощью LLM"»} \\
\end{center}

\vspace{95 pt}

\begin{tabular*}{460pt}{@{\extracolsep{\fill}} l r}
     Студент,\tabularnewline группы 5130201/40003 &  \line(1,0){100} \quad Лебедевич К. О.\\
     & \\
     Руководитель,\tabularnewline преподаватель & \line(1,0){100} \quad Мулюха В.A. \hspace{22 pt} \\
     & \\
     & \\
     & \\
     & \guillemotleft \line(1,0){30} \guillemotright \line(1,0){100} \, 20\,\line(1,0){15} \,г.

\end{tabular*} \\

\vspace{5 pt}

\begin{center}
    Санкт-Петербург, 2025
\end{center}
\newpage

% ----------------------------------------------------------------------
\tableofcontents
\newpage

% ======================================================================
\section{Техническое задание}
% ======================================================================

\textbf{Тема:} Графический калькулятор функций.

\subsection{Цель}

Разработка приложения для построения графиков математических функций одной
переменной в декартовой системе координат. Приложение должно предоставлять
интуитивно понятный интерфейс с программной клавиатурой для ввода выражений
и визуализацией графиков в фиксированном диапазоне.

\subsection{Границы проекта}

\subsubsection*{Включено}

\begin{itemize}
    \item Построение графиков функций одной переменной $y = f(x)$.
    \item Фиксированный диапазон отображения: $x \in [-10^3, 10^3]$,
          $y \in [-10^3, 10^3]$.
    \item Одновременное отображение только одной функции.
    \item Программная клавиатура для ввода математических выражений.
    \item Поддерживаемые математические операции и функции:
    \begin{itemize}
        \item Арифметические операции: $+$, $-$, $*$, $/$, \^{}~(степень).
        \item Тригонометрические функции: $\sin(x)$, $\cos(x)$, $\tan(x)$.
        \item Обратные тригонометрические функции: $\arcsin(x)$,
              $\arccos(x)$, $\arctan(x)$.
        \item Гиперболические функции: $\sinh(x)$, $\cosh(x)$, $\tanh(x)$.
        \item Логарифмические функции: $\log_{10}(x)$, $\ln(x)$.
        \item Степенные функции: $\sqrt{x}$, $x^a$.
        \item Модуль: $|x|$.
        \item Экспонента: $e^x$.
        \item Константы: $\pi$, $e$.
    \end{itemize}
    \item Автоматическая расстановка знака умножения (например,
          $2x$ интерпретируется как $2 \cdot x$).
    \item Обработка унарного минуса.
    \item Проверка корректности введённого выражения.
    \item Отображение осей координат с подписями.
    \item Отображение сетки с шагом 10 единиц.
    \item Отображение текущей позиции курсора в поле ввода.
    \item Статусная строка с информацией о текущем состоянии.
    \item Управление с клавиатуры: Enter (построить график),
          Esc (очистить), Backspace (удалить символ),
          стрелки (перемещение курсора).
\end{itemize}

\subsubsection*{Исключено}

\begin{itemize}
    \item Параметрические функции.
    \item Функции нескольких переменных.
    \item Анимация и динамическое изменение графиков.
    \item Масштабирование и панорамирование (фиксированный масштаб).
    \item Сохранение графиков в файл.
    \item Печать графиков.
    \item Комплексные числа.
\end{itemize}

\subsection{Ограничения}

\subsubsection*{Ограничения на аргументы и значения}

\begin{itemize}
    \item Минимальное значение аргумента: $-10^5$.
    \item Максимальное значение аргумента: $10^5$.
    \item Минимальное отображаемое значение функции: $-10^5$.
    \item Максимальное отображаемое значение функции: $10^5$.
    \item Количество точек для построения графика: 2000 (равномерно
          распределены по диапазону).
\end{itemize}

\subsubsection*{Ограничения на ввод}

\begin{itemize}
    \item Максимальная длина выражения: 256 символов.
    \item Допустимые символы: цифры (0--9), латинские буквы (a--z, A--Z),
          знаки операций ($+, -, *, /, $\^{}$)$, скобки, точка.
    \item Имена функций: \texttt{sin}, \texttt{cos}, \texttt{tan},
          \texttt{asin}, \texttt{acos}, \texttt{atan},
          \texttt{sinh}, \texttt{cosh}, \texttt{tanh},
          \texttt{log}, \texttt{ln}, \texttt{sqrt}, \texttt{abs},
          \texttt{exp}.
    \item Константы: \texttt{pi}, \texttt{e}. Переменная: \texttt{x}
          (только строчная).
\end{itemize}

\subsubsection*{Ограничения вычислителя}

\begin{itemize}
    \item Собственная реализация парсера математических выражений.
    \item Стандартный математический приоритет операций.
    \item Поддержка вложенных функций и скобок.
    \item Обработка ошибок: деление на ноль, логарифм отрицательного
          числа, корень из отрицательного числа.
\end{itemize}

\subsection{Функциональные требования}

\textbf{Интерфейс пользователя} должен содержать строку ввода
с мигающим курсором, программную клавиатуру (цифры, операторы, скобки,
константы $\pi, e$, переменную $x$, функции, управляющие кнопки
$C$, $\leftarrow$, $\to$, $=$, $\Leftarrow$), область построения
с осями, сеткой, графиком (синяя линия) и точкой начала координат,
а также статусную строку.

\textbf{Парсер} должен выполнять лексический и синтаксический анализ,
вычислять значение выражения для заданного $x$, обрабатывать унарный
минус, неявное умножение и константы.

\textbf{Построение} ведётся по 2000 точкам, равномерно распределённым
в диапазоне $x \in [-10^5, 10^5]$. Линии через разрывы функции
пропускаются.

\textbf{Взаимодействие.} Ввод возможен только через программную клавиатуру
(физическая клавиатура блокируется, кроме клавиш Enter/Esc/Backspace/$\leftarrow$/$\to$).
При нажатии на кнопку функции автоматически добавляются скобки и курсор
помещается внутрь.

\subsection{Обработка ошибок}

\begin{itemize}
    \item Пустое выражение --- \textit{«Ошибка: Введите функцию»}.
    \item Некорректное выражение --- \textit{«Ошибка: Некорректное выражение»}.
    \item Деление на ноль, логарифм / корень из отрицательного числа,
          переполнение --- пропуск соответствующих точек графика.
    \item Не удалось построить график --- \textit{«Ошибка: Не удалось
          построить график»}.
\end{itemize}

\subsection{Технологии реализации}

\begin{itemize}
    \item Язык программирования: C++20.
    \item Фреймворк: Qt 6 Widgets.
    \item Сборочная система: CMake.
    \item Компилятор: Clang (Apple LLVM).
    \item Графика: \texttt{QPainter}.
\end{itemize}

\subsection{Отклонения от ТЗ, согласованные с заказчиком}

В процессе разработки заказчик внёс два уточнения, расширяющие исходное ТЗ:
\begin{enumerate}
    \item Пределы по осям ($x_{\min}$, $x_{\max}$, $y_{\min}$, $y_{\max}$)
          и шаг сетки сделаны \emph{регулируемыми}, а не фиксированными.
    \item Добавлено \emph{масштабирование колесом мыши} и
          \emph{панорамирование} перетаскиванием левой кнопкой.
    \item Количество точек графика \emph{вычисляется автоматически}
          (пропорционально ширине области построения и ограничено
          разумным диапазоном) и не задаётся пользователем.
\end{enumerate}

\newpage

% ======================================================================
\section{История промптов}
% ======================================================================

В этом разделе приведена хронологическая последовательность запросов,
переданных языковой модели в ходе разработки. Тексты приводятся
в исходной формулировке пользователя.

\subsection{Промпт 1 --- первичная реализация}

\begin{quote}
\itshape
«Реализуй приложение в папке \texttt{graphicCalculator} в соответствии
с техническим заданием. Но --- в отличии от ТЗ сделай регулируемые
изменяемые пределы и шаг.»
\end{quote}

\textbf{Результат:} модель создала структуру проекта (CMake-сборка,
точка входа, модули \texttt{src/parser/}, \texttt{src/widgets/},
\texttt{MainWindow}), реализовала собственный парсер математических
выражений (лексер, синтаксический анализатор методом рекурсивного
спуска, AST, вычислитель), а также все требуемые виджеты: поле ввода
с блокировкой физической клавиатуры, экранную клавиатуру,
область построения графика и панель настроек с регулируемыми пределами
и шагом сетки.

\subsection{Промпт 2 --- запуск приложения}

\begin{quote}
\itshape
«Как запустить приложение?»
\end{quote}

\textbf{Результат:} модель привела последовательность команд CMake
(\texttt{cmake -B build}, \texttt{cmake --build build -j},
\texttt{open build/GraphicCalculator.app}), указала путь к исполняемому
файлу внутри \texttt{.app}-бандла и привела вариант запуска
без бандла для отладки.

\subsection{Промпт 3 --- интерактивное масштабирование}

\begin{quote}
\itshape
«Добавь масштабирование используя мышку. Также пользователь не должен
вводить число точек --- оно рассчитывается программой.»
\end{quote}

\textbf{Результат:} модель добавила обработчики \texttt{wheelEvent},
\texttt{mousePressEvent}, \texttt{mouseMoveEvent},
\texttt{mouseReleaseEvent} в виджет графика. Колесо мыши реализует
зум вокруг точки курсора (множитель $1.2$ за «деление»),
левая кнопка мыши --- панорамирование. При наведении показывается
всплывающая подсказка с математическими координатами курсора
и значением $f(x)$. Сигнал \texttt{viewChanged} синхронизирует
спинбоксы панели настроек без рекурсии. Поле «Количество точек»
удалено из интерфейса; число выборок вычисляется автоматически
как $\mathrm{clamp}(2 \cdot W_{\text{plot}},\, 500,\, 8000)$.

\subsection{Промпт 4 --- отчёт}

\begin{quote}
\itshape
«Напиши отчёт кодом LaTeX о проделанной работе. Отчёт должен содержать
титульный лист, содержание, ТЗ, историю промптов, описание использованной
языковой модели, заключение и выводы (достоинства и недостатки
реализованной программы), ссылку на GitHub исходного кода.»
\end{quote}

\textbf{Результат:} настоящий документ.

\newpage

% ======================================================================
\section{Описание использованной языковой модели}
% ======================================================================

\subsection{Семейство и версия}

В качестве инструмента разработки использовалась языковая модель
\textbf{Claude Opus 4.7} (идентификатор \texttt{claude-opus-4-7})
производства \textbf{Anthropic} в режиме интерактивного агента
\textit{Claude Code} --- официального CLI-клиента Anthropic для
работы с исходным кодом.

\subsection{Технические характеристики}

\begin{itemize}
    \item Семейство: Claude 4.x (флагманская линейка Anthropic
          на момент проведения работ).
    \item Уровень модели: Opus --- старшая по возможностям модель в семействе.
    \item Дата окончания обучающего корпуса: январь 2026 г.
    \item Поддерживаемая длина контекста --- сотни тысяч токенов;
          в режиме Claude Code контекст автоматически уплотняется
          при приближении к лимиту.
    \item Модальность: текст; в режиме Claude Code модель оснащена
          набором инструментов для работы с файловой системой
          (\texttt{Read}, \texttt{Write}, \texttt{Edit}), запуска
          команд оболочки (\texttt{Bash}), поиска по проекту,
          ведения списка задач, а также рядом специализированных агентов.
\end{itemize}

\subsection{Особенности применения}

В рамках данной работы модель использовалась как:
\begin{enumerate}
    \item Архитектор --- проектирование разбиения на модули
          (парсер, виджеты, главное окно, панель настроек).
    \item Исполнитель --- генерация исходных файлов C++/Qt
          и конфигурации CMake.
    \item Отладчик --- запуск сборки, чтение сообщений компилятора
          и автоматическое исправление ошибок (в частности, конфликта
          типов \texttt{int} и \texttt{qsizetype} при сравнении длины строки).
    \item Технический писатель --- подготовка пользовательской
          документации и настоящего отчёта.
\end{enumerate}

Все правки в коде вносились через инструменты \texttt{Edit}/\texttt{Write},
сборка выполнялась через \texttt{cmake --build}, а корректность запуска
проверялась запуском бинарного файла в режиме \texttt{-platform offscreen},
что подходит для headless-проверки Qt-приложений.

\newpage

% ======================================================================
\section{Архитектура реализованного приложения}
% ======================================================================

Приложение состоит из четырёх логических слоёв.

\begin{enumerate}
    \item \textbf{Парсер выражений} (\texttt{src/parser/}):
    \begin{itemize}
        \item \texttt{Tokenizer} --- лексический анализатор;
              разбивает строку на токены и автоматически вставляет
              знак умножения между смежными «открывающим» и «закрывающим»
              токенами (например, между числом и переменной).
        \item \texttt{Parser} --- рекурсивный спуск с грамматикой
              \emph{expression $\to$ term $\to$ unary $\to$ power $\to$ primary};
              правоассоциативная степень, поддержка вложенных скобок и
              унарного минуса. Возвращает AST в виде \texttt{std::unique\_ptr<AstNode>}.
        \item \texttt{Evaluator} --- рекурсивный обход AST;
              при ошибке домена (\texttt{log} от $\leq 0$, $\sqrt{\cdot}$
              от $<0$, деление на ноль, переполнение, нецелая степень
              от отрицательного основания) возвращает \texttt{NaN},
              что трактуется как «пропустить точку».
    \end{itemize}
    \item \textbf{Виджеты} (\texttt{src/widgets/}):
    \begin{itemize}
        \item \texttt{InputField} --- наследник \texttt{QLineEdit};
              переопределяет \texttt{keyPressEvent} и пропускает
              только Enter, Esc, Backspace и стрелки, эмулируя
              требование «ввод только через программную клавиатуру».
        \item \texttt{KeyboardWidget} --- сетка кнопок, сгруппированных
              по категориям (функции, цифры, операторы, константы,
              управляющие клавиши).
        \item \texttt{GraphWidget} --- отрисовка через \texttt{QPainter}:
              сетка, оси с подписями, кривая графика; обрабатывает
              колесо мыши (зум вокруг точки курсора) и перетаскивание
              (панорама); количество выборок вычисляется по ширине
              области построения.
        \item \texttt{SettingsPanel} --- спинбоксы пределов и шага сетки;
              синхронизируется с интерактивным масштабированием через
              сигнал \texttt{viewChanged}.
    \end{itemize}
    \item \textbf{Главное окно} (\texttt{src/MainWindow.\{h,cpp\}}) ---
          компоновка, состояние парсера, статусная строка
          с индикатором позиции курсора, обработка ошибок и сообщения
          пользователю.
    \item \textbf{Сборка} --- \texttt{CMakeLists.txt} с \texttt{qt\_add\_executable},
          автоматический \texttt{moc}/\texttt{rcc}/\texttt{uic}, поиск Qt
          по пути \texttt{/opt/homebrew/opt/qt} на macOS.
\end{enumerate}

% ======================================================================
\section{Заключение и выводы}
% ======================================================================

\subsection{Итоги}

В ходе работы было разработано полнофункциональное настольное приложение
для построения графиков функций одной переменной, удовлетворяющее
всем требованиям технического задания. Дополнительно к ТЗ реализованы:
\begin{itemize}
    \item регулируемые пределы и шаг сетки;
    \item масштабирование колесом мыши и панорамирование перетаскиванием;
    \item автоматический подбор количества выборок в зависимости от
          размеров области построения;
    \item всплывающая подсказка с координатами курсора и значением
          функции в точке.
\end{itemize}

\subsection{Достоинства}

\begin{itemize}
    \item \textbf{Собственный парсер} реализован «с нуля»: лексер,
          синтаксический анализатор, AST и вычислитель --- без внешних
          зависимостей кроме Qt. Это удовлетворяет требованию ТЗ
          и облегчает расширение языка выражений.
    \item \textbf{Корректная обработка ошибок:} сообщения на русском
          языке, отдельные сообщения для пустого выражения, синтаксической
          ошибки и пустого графика; разрывы и недопустимые значения
          вызывают пропуск точек, а не аварийное завершение.
    \item \textbf{Расширенная интерактивность:} зум вокруг точки курсора,
          панорама, тултип со значениями --- решения, естественные
          для современного UX.
    \item \textbf{Автоматическая адаптация плотности выборки:} избавляет
          пользователя от выбора числа точек и обеспечивает гладкость
          при любом размере окна.
    \item \textbf{Чистая модульная архитектура:} один класс ---
          одна ответственность, что облегчает сопровождение.
    \item \textbf{Кроссплатформенность} обеспечивается Qt 6 и CMake.
\end{itemize}

\subsection{Недостатки}

\begin{itemize}
    \item Одновременно отображается только \emph{одна} функция;
          сравнение нескольких графиков невозможно.
    \item Шаг сетки не подстраивается автоматически при сильном
          изменении масштаба --- при значительном зуме пользователь
          должен вручную поправить шаг.
    \item Отсутствует история введённых выражений и возможность
          сохранения / загрузки сессии.
    \item Численная обработка асимптот эвристическая: разрыв определяется
          по «большому скачку» $|\Delta y|$; для редких функций возможны
          артефакты в виде разрыва там, где его нет, или короткой
          соединительной линии.
    \item Нет аналитических возможностей: производных, нулей функции,
          точек пересечения, экстремумов.
    \item Шрифт \texttt{Menlo}, использованный в стилях, доступен
          из коробки только на macOS; на других платформах будет
          выбран замещающий шрифт.
    \item Жёсткие границы плоскости ($\pm 10^5$) могут оказаться
          узкими для отдельных классов задач.
\end{itemize}

\subsection{Выводы}

Применение крупной языковой модели Claude Opus 4.7 в режиме агента
Claude Code позволило за четыре итеративных запроса довести проект
от пустой директории до работающего приложения, удовлетворяющего ТЗ
и расширенного интерактивными возможностями. Модель самостоятельно
проектировала модули, генерировала исходный код, исправляла ошибки
компиляции и адаптировала архитектуру при изменении требований
заказчика (переход от фиксированных параметров к регулируемым,
а затем --- к интерактивному управлению мышью).

Такой подход существенно сокращает время на реализацию типовых
desktop-приложений и пригоден как минимум для прототипирования
и учебных проектов. Главным ограничением остаётся необходимость
содержательной проверки сгенерированного кода человеком ---
особенно в местах, где автоматический тест не выявляет логических
дефектов (например, эвристика разрыва функции).

\newpage

% ======================================================================
\section{Ссылка на исходный код}
% ======================================================================

Исходный код приложения и настоящий отчёт размещены в репозитории
на GitHub:

\begin{center}
\large
\href{https://github.com/Lenebik/GraphicCalculator}{\texttt{https://github.com/Lenebik/GraphicCalculator}}
\end{center}

\noindent
Структура репозитория:
\begin{itemize}
    \item \texttt{graphicCalculator/} --- исходный код приложения
          (C++20, Qt 6 Widgets, CMake);
    \item \texttt{graphicCalculator/src/parser/} --- модуль разбора
          и вычисления математических выражений;
    \item \texttt{graphicCalculator/src/widgets/} --- виджеты UI:
          поле ввода, экранная клавиатура, область графика,
          панель настроек;
    \item \texttt{graphicCalculator/CMakeLists.txt} --- конфигурация
          сборки;
    \item \texttt{report.tex} --- исходник настоящего отчёта.
\end{itemize}

\end{document}
